
\documentclass[11pt,a4paper]{article}
\usepackage[spanish]{babel}
\usepackage[utf8]{inputenc}
\usepackage[T1]{fontenc}
\usepackage{amsmath,amssymb}
\usepackage{siunitx}
\usepackage{geometry}
\geometry{margin=2.5cm}

\begin{document}

\title{Línea tanque--proceso: planteo del modelo adiabático con \texorpdfstring{$P_0$}{P0} desconocida}
\author{}
\date{}
\maketitle

\section*{1. Subproblema que se modela}

Consideramos sólo la \textbf{línea tanque--proceso}:

\begin{itemize}
  \item Reservorio 0: el tanque, a presión desconocida $P_0$ y temperatura aproximadamente constante
  \[
    T_0 \simeq 40^\circ\text{C} \approx \SI{313}{K}.
  \]
  \item Reservorio 3: el proceso, que opera a presión
  \[
    P_3 = \SI{2}{bar(g)} = \SI{3}{bar(abs)}.
  \]
  \item Conducción: tubería de acero Sch 40, 1/2'' nominal, longitud equivalente
  \[
    L'_e = \SI{30}{m},\qquad D \simeq \SI{0.0158}{m}.
  \]
  \item Gas: aire, tratado como gas ideal:
  \[
    M = \SI{29}{g/mol},\quad
    \gamma = \frac{c_p}{c_v} = 1.4,\quad
    R = \frac{R_u}{M} \approx \frac{8.314}{0.029}\simeq 2.87\times10^2\ \text{J/(kg\,K)}.
  \]
  \item Factor de fricción Darcy tabulado (Crane) para 1/2'' Sch 40, régimen turbulento completo:
  \[
    f_T \simeq 0.026.
  \]
\end{itemize}

Condición de diseño en régimen permanente:
\[
\dot m = \dot m_{\text{proc}} = \SI{2}{kg/min} =
\frac{2}{60}\ \text{kg/s} \approx 3.33\times10^{-2}\ \text{kg/s}.
\]

El objetivo de esta parte es expresar el \textbf{sistema de ecuaciones no lineales reducido} cuyas incógnitas son $(r,z,F)$ (o, equivalentemente, $P_0$), para este flujo adiabático con fricción.

\section*{2. Hipótesis físicas y elección del modelo}

Para la línea se adoptan las siguientes hipótesis:

\begin{enumerate}
  \item \textbf{Flujo estacionario} (aproximación cuasi-permanente en régimen).
  \item \textbf{Gas ideal}, propiedades constantes $(R,\gamma)$.
  \item \textbf{Conducción horizontal, sección constante}, sin trabajo de eje.
  \item \textbf{Modelo adiabático con fricción} (tipo Fanno + reservorio) como modelo ideal:

  El tanque está a $T_0\simeq 40^\circ\text{C}$ y el ambiente a $15^\circ\text{C}$.
  El gas se enfría a lo largo de la línea ($Q<0$).
  Del análisis comparativo de modelos (S4):
  \[
  \left(\frac{w}{A}\right)_{\text{iso}}
  < \left(\frac{w}{A}\right)_{\text{adiab}}
  < \left(\frac{w}{A}\right)_{\text{real, enfriado}},
  \]
  es decir, el modelo adiabático da un \textbf{flujo menor} que el real; se usa como
  modelo conservador que da la \textbf{mínima capacidad} de la línea para una
  dada caída de presión.
\end{enumerate}

En resumen, se resuelve la línea completa 0--3 como un flujo adiabático desde un reservorio 0 hasta un reservorio 3, subsónico, con $\dot m$ dado y $P_0$ (por ende $F=P_3/P_0$) desconocido.

\section*{3. Definición de variables adimensionales}

Se reutiliza la misma reducción que para ``flujo adiabático desde reservorio, subsónico''.

Volumen específico en el reservorio 0:
\[
v_0 = \frac{R T_0}{P_0}.
\]

Variables adimensionales:
\[
r = \frac{v_0}{v_{1'}},\qquad
z = \frac{v_0}{v_2},
\]
donde $1'$ es la sección de entrada efectiva del ducto y $2$ una sección
interna, que se toma coincidente con las condiciones del proceso, es decir
$P_2 = P_3$.

Relación de presiones (reservorio 0 $\to$ proceso):
\[
F = \frac{P_3}{P_0}.
\]

Parámetro de flujo:
\[
X = \left(\frac{\dot m}{A}\right)^2 \frac{v_0}{P_0},
\]
con
\[
A = \frac{\pi D^2}{4}
\]
el área interna de la tubería.

Además:
\[
C = \frac{2\gamma}{\gamma-1}.
\]
Para aire ($\gamma=1.4$), $C=7$.

\section*{4. Sistema reducido de ecuaciones (forma estructural)}

Del desarrollo adiabático con fricción desde reservorio, las ecuaciones reducidas son:

\begin{enumerate}
  \item Relación reserva--entrada--salida:
  \[
  r^2 - r^{\gamma+1} - z^2 + Fz = 0.
  \]

  \item Definición reducida de $X$ (energía + estado del gas):
  \[
  X = C\,(z^2 - Fz).
  \]

  \item Ecuación de Fanno reducida entre $1'$ y $2$:
  \[
  \frac{r^2 - z^2}{X}
  - \frac{\gamma+1}{\gamma}\ln\left(\frac{r}{z}\right)
  = f_T\,\frac{L'_e}{D}.
  \]
\end{enumerate}

Hasta aquí es el mismo sistema que en el problema p.2, pero allí:

\begin{itemize}
  \item $\dot m$ (y, por tanto, $X$) era incógnita,
  \item $F = P_2/P_0$ era dato.
\end{itemize}

Ahora es al revés:

\begin{itemize}
  \item $\dot m$ es \textbf{dato},
  \item $F$ (y, en consecuencia, $P_0$) es \textbf{incógnita}.
\end{itemize}

\section*{5. Cierre del sistema incorporando \texorpdfstring{$\dot m$}{m punto} dado}

Usamos la definición general de $X$:
\[
X = \left(\frac{\dot m}{A}\right)^2 \frac{v_0}{P_0}
  = \left(\frac{\dot m}{A}\right)^2 \frac{R T_0}{P_0^2}.
\]

Pero $P_0$ se relaciona con $F$ mediante
\[
F = \frac{P_3}{P_0}
\quad\Rightarrow\quad
P_0 = \frac{P_3}{F}.
\]

Entonces:
\[
X = \left(\frac{\dot m}{A}\right)^2
    \frac{R T_0}{\left(\dfrac{P_3}{F}\right)^2}
  = \left(\frac{\dot m}{A}\right)^2
    \frac{R T_0\,F^2}{P_3^2}.
\]

Definimos la constante conocida
\[
K = \left(\frac{\dot m}{A}\right)^2
    \frac{R T_0}{P_3^2},
\]
de modo que
\[
X(F) = K\,F^2.
\]

Del modelo adiabático también tenemos
\[
X = C\,(z^2 - Fz).
\]

Igualando ambas expresiones de $X$, se obtiene una ecuación adicional que liga $z$ y $F$:
\[
C\,(z^2 - Fz) = K\,F^2.
\]

\section*{6. Sistema final en incógnitas \texorpdfstring{$(r,z,F)$}{(r,z,F)}}

El problema de la línea tanque--proceso, en el modelo adiabático reducido con $\dot m$ dado, queda como el sistema no lineal:

\[
\begin{cases}
E_1(r,z,F) = r^2 - r^{\gamma+1} - z^2 + Fz = 0,\\[6pt]
E_2(r,z,F) =
\dfrac{r^2 - z^2}{C\,(z^2 - Fz)}
- \dfrac{\gamma+1}{\gamma}\ln\left(\dfrac{r}{z}\right)
- f_T\dfrac{L'_e}{D} = 0,\\[10pt]
E_3(z,F) = C\,(z^2 - Fz) - K F^2 = 0,
\end{cases}
\]
donde:
\[
\gamma = 1.4,\quad
C = \frac{2\gamma}{\gamma-1} = 7,\quad
f_T = 0.026,\quad
L'_e = \SI{30}{m},\quad
D \simeq \SI{0.0158}{m},
\]
\[
A = \frac{\pi D^2}{4},\quad
T_0 \simeq \SI{313}{K},\quad
P_3 = \SI{3}{bar(abs)} = 3\times10^5\ \text{Pa},
\]
\[
\dot m = \frac{2}{60}\ \text{kg/s},\quad
R \approx 2.87\times10^2\ \text{J/(kg\,K)},
\]
\[
K = \left(\frac{\dot m}{A}\right)^2
      \frac{R T_0}{P_3^2}.
\]

Una vez resuelto este sistema se obtienen $(r,z,F)$ y se recupera
\[
P_0 = \frac{P_3}{F}.
\]

Este $P_0$ es el \textbf{mínimo} que, en el modelo adiabático (conservador aquí), permite que la línea entregue $\dot m = 2\ \text{kg/min}$.

\end{document}
